\section{Squeeze-and-Excitation Blocks\label{sec:se-blocks}}

A Squeeze-and-Excitation block is a computational
unit which can be built upon a transformation $\FF_{tr}$ mapping an input $\XX \in \RR^{H' \times W' \times C'}$ to feature maps $\UU \in \RR^{H \times W \times C}$.  In the notation that follows we take $\FF_{tr}$ to be a convolutional operator and use $\VV= [\vv_1, \vv_2, \dots, \vv_{C}]$ to denote the learned set of filter kernels, where $\vv_c$ refers to the parameters of the \mbox{$c$-th} filter.  We can then write the outputs as $\UU = [\uu_1, \uu_2, \dots, \uu_{C}]$, where

\begin{equation}\label{simple_neuron}
\uu_c = \vv_c \ast \XX = \sum_{s=1}^{C'}\vv^s_c \ast \xx^s.
\end{equation}
Here $\ast$ denotes convolution, $\vv_c = [\vv^1_c, \vv^2_c, \dots, \vv^{C'}_c]$, $\XX = [\xx^1, \xx^2, \dots, \xx^{C'}]$ and \mbox{$\uu_c \in \RR^{H \times W}$}. $\vv^s_c$ is a $2$D spatial kernel representing a single channel of $\vv_c$ that acts on the corresponding channel of $\XX$. To simplify the notation, bias terms are omitted. Since the output is produced by a summation through all channels, channel dependencies are implicitly embedded in $\vv_c$, but are entangled with the local spatial correlation captured by the filters.  The channel relationships modelled by convolution are inherently implicit and local (except the ones at top-most layers). We expect the learning of convolutional features to be enhanced by explicitly modelling channel interdependencies, so that the network is able to increase its sensitivity to informative features which can be exploited by subsequent transformations. Consequently, we would like to provide it with access to global information and recalibrate filter responses in two steps, \textit{squeeze} and \textit{excitation}, before they are fed into the next transformation. A diagram illustrating the structure of an SE block is shown in Fig.~\ref{fig:splash}.


\subsection{Squeeze: Global Information Embedding}
In order to tackle the issue of exploiting channel dependencies, we first consider the signal to each channel in the output features.  Each of the learned filters operates with a local receptive field and consequently each unit of the transformation output $\UU$ is unable to exploit contextual information outside of this region.

To mitigate this problem, we propose to {\it squeeze} global spatial information into a channel descriptor. This is achieved by using global average pooling to generate channel-wise statistics. Formally, a statistic $\zz \in \RR^{C}$ is generated by shrinking $\UU$ through its spatial dimensions $H \times W$, such that the $c$-th element of $\zz$ is calculated by:
%\vspace{-1mm}
\begin{equation}\label{squeeze}
z_c = \FF_{sq}(\uu_c) = \frac{1}{H\times W}\sum_{i=1}^{H} \sum_{j=1}^{W} u_c(i,j).
\end{equation}

{\it Discussion.} The output of the transformation $\UU$ can be interpreted as a collection of the local descriptors whose statistics are expressive for the whole image. Exploiting such information is prevalent in prior feature engineering work \cite{yang_cvpr2009spm, fisher_ijcv2013fv, shen_tip2015}. We opt for the simplest aggregation technique, global average pooling, noting that more sophisticated strategies could be employed here as well.

%-------------------------------------------------------------------------
\subsection{Excitation: Adaptive Recalibration} \label{subsec:adaptive-recal}
To make use of the information aggregated in the \textit{squeeze} operation, we follow it with a second operation which aims to fully capture channel-wise dependencies.  To fulfil this objective, the function must meet two criteria: first, it must be flexible (in particular, it must be capable of learning a nonlinear interaction between channels) and second, it must learn a non-mutually-exclusive relationship since we would like to ensure that multiple channels are allowed to be emphasised (rather than enforcing a one-hot activation). To meet these criteria, we opt to employ a simple gating mechanism with a sigmoid activation:

\begin{equation}\label{excitation}
\mathbf{s} = \mathbf{F}_{ex}(\mathbf{z}, \mathbf{W}) = \sigma(g(\mathbf{z}, \mathbf{W})) = \sigma(\WW_2\delta(\WW_1\mathbf{z})),
\end{equation}
where $\delta$ refers to the ReLU \cite{nair_icml2010relu} function, $\WW_1 \in \RR^{\frac{C}{r} \times C}$ and $\WW_2 \in \RR^{C \times \frac{C}{r}}$. To limit model complexity and aid generalisation, we parameterise the gating mechanism by forming a bottleneck with two fully-connected (FC) layers around the non-linearity, i.e. a dimensionality-reduction layer with reduction ratio $r$ (this parameter choice is discussed in Section~\ref{subsec:reduction}), a ReLU and then a dimensionality-increasing layer returning to the channel dimension of the transformation output $\UU$. The final output of the block is obtained by rescaling $\UU$ with the activations $\mathbf{s}$:

\begin{equation}\label{scale}
\widetilde{\mathbf{x}}_c = \mathbf{F}_{scale}(\mathbf{u}_c, s_c) = s_c\,\mathbf{u}_c,
\end{equation}
where $\widetilde{\XX} = [\widetilde{\xx}_1, \widetilde{\xx}_2, \dots, \widetilde{\xx}_{C}]$ and $\FF_{scale}(\uu_c, s_c)$ refers to channel-wise multiplication between the scalar $s_c$ and the feature map $\uu_c \in \RR^{H \times W}$.

{\it Discussion.}  The excitation operator maps the input-specific descriptor $\mathbf{z}$ to a set of channel weights. In this regard, SE blocks intrinsically introduce dynamics conditioned on the input, which can be regarded as a self-attention function on channels whose relationships are not confined to the local receptive field the convolutional filters are responsive to.
\subsection{Instantiations}
\label{se_inception_resnet}

\begin{figure}[t]
\begin{center}
\includegraphics[height=65mm]{figures/module_seinception.pdf}
\end{center}
\vspace{-2em}
\caption{The schema of the original Inception module (left) and the SE-Inception module (right).}
\label{fig:seinception_module}
\end{figure}

The SE block can be integrated into standard architectures such as VGGNet \cite{simonyan_iclr2015vgg} by insertion after the non-linearity following each convolution. Moreover, the flexibility of the SE block means that it can be directly applied to transformations beyond standard convolutions. To illustrate this point, we develop SENets by incorporating SE blocks into several examples of more complex architectures, described next. 

We first consider the construction of SE blocks for Inception networks \cite{szegedy_cvpr2015googlenet}.  Here, we simply take the transformation $\FF_{tr}$ to be an entire Inception module (see Fig.~\ref{fig:seinception_module}) and by making this change for each such module in the architecture, we obtain an \textit{SE-Inception} network. SE blocks can also be used directly with residual networks (Fig.~\ref{fig:seresnet_module} depicts the schema of an \textit{SE-ResNet} module). Here, the SE block transformation $\FF_{tr}$ is taken to be the non-identity branch of a residual module. \textit{Squeeze} and \textit{Excitation} both act before summation with the identity branch. Further variants that integrate SE blocks with ResNeXt \cite{xie_cvpr2017resnext}, Inception-ResNet \cite{szegedy_iclrw2016inceptionv4}, MobileNet \cite{howard_2017mobilenets} and ShuffleNet \cite{zhang_cvpr2018shufflenet} can be constructed by following similar schemes.  For concrete examples of SENet architectures, a detailed description of \mbox{SE-ResNet-50} and \mbox{SE-ResNeXt-50} is given in Table~\ref{model_structure}.  

One consequence of the flexible nature of the SE block is that there are several viable ways in which it could be integrated into these architectures.  Therefore, to assess sensitivity to the integration strategy used to incorporate SE blocks into a network architecture, we also provide ablation experiments exploring different designs for block inclusion in Section~\ref{sec:integration}.


\begin{figure}
\begin{center}
\includegraphics[height=65mm]{figures/module_seresnet.pdf}
\end{center}
\vspace{-2em}
\caption{The schema of the original Residual module (left) and the SE-ResNet module (right).}
\label{fig:seresnet_module}
\end{figure}
%----------------------------------------------------------------------------